
\section{What sequential transport theorems require}
Nested and adapted transport distances address the temporal issue by
restricting couplings so that matching paths cannot depend on information
that is unavailable at the relevant date. Pflug and Pichler's stability
argument uses convex decision sets, convex loss in decisions and uniform
regularity in uncertainty \cite[Definition~5.1, Theorem~6.1 and proof]{litclosepflug}.
A projected action is formed by conditional averaging. If the only permitted
issuance sizes are 0 and 1, an average of $1/2$ is infeasible. The theorem
therefore does not automatically justify merging discrete issuance choices.

Backhoff-Veraguas and coauthors obtain related stability results for financial
optimization with bounded predictable strategies and regular payoffs
\cite[Theorem~1.5; Lemma~3.7; Theorem~3.10]{litcloseadapted}. Their weak transfer
permits a different strategy in the other model. Reusing the same strategy
needs additional regularity. Their semimartingale distance also has a
different cost functional from a simple sum of discrete path deviations; the
paper explains why time refinement matters. Its examples of discontinuous
strategies and exact superhedging delimit the guarantee. These are useful
reasons to formulate the company decision problem precisely, rather than
reasons to attach one distance to every accounting query.
